\section{Original well-founded canonical-segment prerequisite}

\begin{theorem}[The maximum drops along canonical segments; well-founded values resolve]
\label{hir:step-one-well-founded}
\uses{hir:step-one-bridge,hir:centre-function,hir:resolution-driver,hir:invariant-driver,hir:canonical-segments,bcalc:transform-off-centre,bcalc:inv-domain-stable}
\source{wlodarczyk05-hironaka,bierstone-milman08-functoriality,kollar07-resolution}
\lean{Hironaka.StepOneWellFounded.StepDrop,Hironaka.StepOneWellFounded.lt_of_mem_maxLocus,Hironaka.StepOneWellFounded.lt_of_isCanonical_append,Hironaka.StepOneWellFounded.lt_of_isCanonical_step,Hironaka.StepOneWellFounded.proj_mem_cosupp_step}\upstreamlean{TranscendentalNakamaye.Pending.StepOneWellFounded.stepDrop_of_resolutionInvariant,TranscendentalNakamaye.Pending.StepOneWellFounded.exists_isGeoResolution_of_wellFoundedLT}\lean{Hironaka.StepOneWellFounded.stepDrop_stepOneInvM}\upstreamlean{TranscendentalNakamaye.Pending.StepOneWellFounded.stepOneInvM_lt_of_isCanonical_append,TranscendentalNakamaye.Pending.StepOneWellFounded.isAdmissibleCentre_maxLocus_stepOneInvM,TranscendentalNakamaye.Pending.StepOneWellFounded.exists_isGeoResolution_stepOneInvM}
\mixedleanok
Let \(\underline{\mathcal I}\) be a marked ideal on \(X\) and \(\operatorname{inv}\) a function on the towers over \(X\) such that after every step of a canonical segment
(the centre \(C\) empty or the maximum locus) every point \(z\) of the new cosupport lies over the cosupport, with \(\operatorname{inv}(z)<\operatorname{inv}(\pi z)\) if
\(\pi z\in C\) and \(\operatorname{inv}(z)=\operatorname{inv}(\pi z)\) otherwise.
\begin{enumerate}
\item If \(C\neq\emptyset\) is the maximum locus after a canonical segment \(T\), with value \(\lambda\), then along every canonical segment extending \(T\) by the blowup of
\(C\) every value on the cosupport is \(<\lambda\).
\item A resolution invariant (Definition~\ref{hir:centre-function}) has this property along its canonical segments; if its values are well-ordered, it resolves every
marked ideal of its domain over a compact base by a canonical segment without empty centres, without using property (8).
\item The masked Step 1 value of a marked ideal of maximal order in the domain has this property (with the maximum loci on \(\operatorname{cosupp}\underline{\mathcal J}_T\)),
its maximum locus is an admissible centre, and if the values \(\Lambda\) of the lower invariant are well-ordered, it resolves \(\underline{\mathcal J}\) over a compact base
by a canonical segment without empty centres.
\end{enumerate}
\end{theorem}
\begin{proof}
\mixedleanok
(1) By induction along the extension. After the blowup of \(C\): a point over \(C\) has a value \(<\operatorname{inv}(\pi z)=\lambda\); a point off \(C\) keeps the value of
\(\pi z\), a point of the cosupport outside the maximum locus, so \(<\lambda\). At a later step every value is \(<\) or \(=\) the value of its image, which is \(<\lambda\) by
induction. The points of the cosupport lie over the cosupport: over the centre since the centre is admissible, off it since the blowup is an isomorphism there
(Lemma~\ref{bcalc:transform-off-centre}).
(2) Properties (3) and (4) of Definition~\ref{hir:centre-function} give the hypothesis, canonical segments being multiple test blow-ups
(Theorem~\ref{hir:invariant-driver}). With (1) for one step, the hypotheses of the well-founded form of Theorem~\ref{hir:resolution-driver} hold: closed cosupports
(Lemma~\ref{bcalc:inv-domain-stable}), (2), admissible maximum loci with empty interior (Theorem~\ref{hir:invariant-driver}), and the drop.
(3) The same with Corollary~\ref{hir:step-one-bridge} in place of the fields; \(\mathbf N\times(\Lambda\cup\{\infty\})\), ordered lexicographically, is well-ordered
when \(\Lambda\) is.
\end{proof}

\begin{lemma}[Cosupports of sums, products and powers]
\label{hir:marked-ideal-operations}
\uses{hir:cosupport-basic,hir:order-basic,hir:order-power,hir:derivative-leibniz}
\source{wlodarczyk05-hironaka}
\lean{Hironaka.IdealOrderArithmetic.vanishesToOrderOn_sum_mul,Hironaka.IdealOrderArithmetic.IsInSpanAt.vanishesToOrderOn,Hironaka.IdealOrderArithmetic.cosupport_sumIdeal,Hironaka.IdealOrderArithmetic.inter_cosupport_subset_cosupport_mul,Hironaka.IdealOrderArithmetic.cosupport_subset_cosupport_idealPow,Hironaka.OrderMultiplicativity.cosupport_idealPow_mul_eq}\upstreamlean{TranscendentalNakamaye.Pending.OrderMultiplicativity.cosupport_sumIdeal_idealPow}
\mixedleanok
Let \(\mathfrak a,\mathfrak b\) be ideals of holomorphic germs and \(N\subset M\). Then
\(\operatorname{cosupp}(\mathfrak a+\mathfrak b|_N,\mu)=\operatorname{cosupp}(\mathfrak a|_N,\mu)\cap\operatorname{cosupp}(\mathfrak b|_N,\mu)\),
\(\operatorname{cosupp}(\mathfrak a|_N,\mu)\cap\operatorname{cosupp}(\mathfrak b|_N,\nu)\subset\operatorname{cosupp}(\mathfrak a\mathfrak b|_N,\mu+\nu)\) and
\(\operatorname{cosupp}(\mathfrak a|_N,\mu)\subset\operatorname{cosupp}(\mathfrak a^\nu|_N,\mu\nu)\). For \(N=M\) and \(\nu\ge1\),
\(\operatorname{cosupp}(\mathfrak a^\nu,\mu\nu)=\operatorname{cosupp}(\mathfrak a,\mu)\); hence for \(\mu,\nu\ge1\) the sum of marked ideals
\((\mathfrak a,\mu)+(\mathfrak b,\nu)=(\mathfrak a^\nu+\mathfrak b^\mu,\mu\nu)\) has cosupport
\(\operatorname{cosupp}(\mathfrak a,\mu)\cap\operatorname{cosupp}(\mathfrak b,\nu)\).
\end{lemma}
\begin{proof}
\mixedleanok
Germs vanishing to order \(\mu\) along \(N\) at \(x\) are closed under holomorphic combinations, and products add orders
(Lemma~\ref{hir:order-basic}); this gives the first three statements (the first in both directions because
\(\mathfrak a,\mathfrak b\subset\mathfrak a+\mathfrak b\)), the third by induction on \(\nu\). For the reverse inclusion of the fourth, if
\(x\in\operatorname{cosupp}(\mathfrak a^\nu,\mu\nu)\) and \(g\in\mathfrak a_x\), then \(g^\nu\in(\mathfrak a^\nu)_x\subset\mathfrak m_x^{\nu\mu}\), so \(g\in\mathfrak m_x^\mu\) by
Lemma~\ref{hir:order-power}. The last statement combines the first and the fourth.
\end{proof}

\begin{lemma}[Controlled transforms of products, powers and sums]
\label{hir:controlled-operations}
\uses{bcalc:controlled-transform,bcalc:exceptional-divisor}
\source{wlodarczyk05-hironaka}
\lean{Hironaka.MarkedIdealOperationsTransform.mul_le_mul',Hironaka.MarkedIdealOperationsTransform.idealPow_mono,Hironaka.MarkedIdealOperationsTransform.idealPow_idealPow,Hironaka.MarkedIdealOperationsTransform.mul_sumIdeal,Hironaka.MarkedIdealOperationsTransform.pullback_sumIdeal,Hironaka.MarkedIdealOperationsTransform.isCoherent_D,Hironaka.MarkedIdealOperationsTransform.controlledTransform_mul,Hironaka.MarkedIdealOperationsTransform.controlledTransform_idealPow,Hironaka.MarkedIdealOperationsTransform.controlledTransform_sumIdeal}\upstreamlean{TranscendentalNakamaye.Pending.MarkedIdealOperationsTransform.controlledTransform_sumIdeal_idealPow}
\mixedleanok
Let \(\sigma\) be a blowup along \(C\) and \(\mathfrak a\subset\mathcal I_C^\mu\), \(\mathfrak b\subset\mathcal I_C^\nu\) coherent. Then
\(\sigma^{\mathrm c}(\mathfrak a\mathfrak b,\mu+\nu)=\sigma^{\mathrm c}(\mathfrak a,\mu)\,\sigma^{\mathrm c}(\mathfrak b,\nu)\), \(\sigma^{\mathrm c}(\mathfrak a^k,\mu k)=\sigma^{\mathrm c}(\mathfrak a,\mu)^k\),
\(\sigma^{\mathrm c}(\mathfrak a+\mathfrak b,\mu)=\sigma^{\mathrm c}(\mathfrak a,\mu)+\sigma^{\mathrm c}(\mathfrak b,\mu)\) if also \(\mathfrak b\subset\mathcal I_C^\mu\), and
\(\sigma^{\mathrm c}(\mathfrak a^\nu+\mathfrak b^\mu,\mu\nu)=\sigma^{\mathrm c}(\mathfrak a,\mu)^\nu+\sigma^{\mathrm c}(\mathfrak b,\nu)^\mu\).
\end{lemma}
\begin{proof}
\mixedleanok
Pullback commutes with products and sums of ideals and \(\mathfrak l(\mathfrak a+\mathfrak b)=\mathfrak l\mathfrak a+\mathfrak l\mathfrak b\); so, with \(\mathcal O(-\mu D)\) the
\(\mu\)-th power of the exceptional ideal, \(\sigma^*(\mathfrak a\mathfrak b)=\mathcal O(-(\mu+\nu)D)\,\sigma^{\mathrm c}(\mathfrak a,\mu)\sigma^{\mathrm c}(\mathfrak b,\nu)\),
\(\sigma^*(\mathfrak a^k)=\mathcal O(-\mu kD)\,\sigma^{\mathrm c}(\mathfrak a,\mu)^k\) and
\(\sigma^*(\mathfrak a+\mathfrak b)=\mathcal O(-\mu D)(\sigma^{\mathrm c}(\mathfrak a,\mu)+\sigma^{\mathrm c}(\mathfrak b,\mu))\), and the controlled transform is the unique coherent
ideal with this property (Lemma~\ref{bcalc:controlled-transform}). For the last statement, \(\mathfrak a^\nu\subset\mathcal I_C^{\mu\nu}\) and
\(\mathfrak b^\mu\subset\mathcal I_C^{\mu\nu}\); apply the sum rule and then the power rule.
\end{proof}

\begin{lemma}[Deleting the members that contain every centre]
\label{hir:drop-containing-members}
\uses{bcalc:nc-comparison,bcalc:sub-divisor,bcalc:nc-comparison-transform,hir:drop-members,hir:snc,mc:strict-along,hir:geo-transform-basic,hir:test-tower,hir:wellformed-tower}
\source{wlodarczyk05-hironaka}
\lean{Hironaka.DropMembers.isGeoTestTower_withE_iff_of_centresIn_of_snc,Hironaka.DropMembers.isGeoTestTower_withE_iff_of_centresIn}
\leanok
Let \(\underline{\mathcal I}=(M,N,E,\mathfrak a,d)\) be well formed, \(E'\) a list of members of \(E\), and \(S\) a set contained in every member of \(E\) missing from
\(E'\) (for instance \(S=N\cap\bigcap_{H\in A}H\) and \(E'=E_S\) the sequence \(E\) without the members of a set \(A\)). For towers whose centres lie in the
strict transforms \(S_T\) of \(S\): \(T\) is a multiple test blow-up of \(\underline{\mathcal I}\) if and only if it is one of
\(\underline{\mathcal I}'=(M,N,E',\mathfrak a,d)\), and along it the two transforms have the same submanifold, ideal and cosupport (Lemma~\ref{hir:drop-members});
their sequences differ by members containing \(S_T\).
\end{lemma}
\begin{proof}
\leanok
Induction on the tower. The members of the transformed \(E'\) are members of the transformed \(E\) (strict transforms of members of \(E'\), and the exceptional
divisor, which is in both). A member of the transformed \(E\) missing from the transformed \(E'\) is the strict transform of a member \(H\) missing from \(E'\),
and \(S_T\subset H_T\) by induction; strict transforms are monotone, so it contains \(S_{T'}\) after the next blowup. A centre \(C\subset S_T\) therefore
lies in every member missing from the transformed \(E'\), and, the transform of \(\underline{\mathcal I}\) being well formed along a multiple test blow-up
(Lemma~\ref{hir:wellformed-tower}), its divisor has simple normal crossings and the normal-crossings comparison (Theorem~\ref{bcalc:nc-comparison}) shows that \(C\)
has only normal crossings with one divisor if and only if with the other. Submanifold, ideal and cosupport agree by Lemma~\ref{hir:drop-members}, so \(C\) is
admissible for one transform if and only if for the other. Well-formedness is used only for the simple normal crossings of the transformed \(E\) along
the multiple test blow-ups of \(\underline{\mathcal I}\); the lemma holds whenever these hold.
\end{proof}
